% ============================================================
% Kapitel 1: Einleitung
% ============================================================

\section{Einleitung}
\label{sec:einleitung}

Lorem ipsum dolor sit amet, consectetur adipiscing elit, sed do eiusmod tempor
incididunt ut labore et dolore magna aliqua. Ut enim ad minim veniam, quis nostrud
exercitation ullamco laboris nisi ut aliquip ex ea commodo consequat. Duis aute
irure dolor in reprehenderit in voluptate velit esse cillum dolore eu fugat nulla
pariatur.\autocite[Vgl.][S. 1]{muster:buch:2023}

Lorem ipsum dolor sit amet, consectetur adipiscing elit, sed do eiusmod tempor
incididunt ut labore et dolore magna aliqua. Ut enim ad minim veniam, quis nostrud
exercitation ullamco laboris nisi ut aliquip ex ea commodo consequat.\autocite[Vgl.][S. 100]{muster:artikel:2022}

% ── Beispiel: Blockzitat ──────────────────────────────────────────────────────
% Für wörtliche Zitate ab ca. 3 Zeilen: eingerückt, kleinere Schrift, einzeilig.
% \autocite[Vgl.][S. 42]{schluessel}          → Vgl.-Fußnotenzitat
% \autocite[][S. 42]{schluessel}               → direktes Fußnotenzitat
% \autocites[Vgl.][S. 1]{q1}[Vgl.][S. 5]{q2} → mehrere Quellen gleichzeitig
% \cite[S. 42]{schluessel}                     → inline (z.B. in Bildunterschriften)

\begin{quotation}
	„Zitat 1: Lorem ipsum dolor sit amet, consectetur adipiscing elit, sed do
	eiusmod tempor incididunt ut labore et dolore magna aliqua. Ut enim ad minim
	veniam, quis nostrud exercitation ullamco laboris nisi ut aliquip."\autocite[][S. 55]{muster:kapitel:2021}
\end{quotation}

Lorem ipsum dolor sit amet, consectetur adipiscing elit, sed do eiusmod tempor
incididunt ut labore et dolore magna aliqua.

% ── Beispiel: Aufzählung (ungeordnet) ────────────────────────────────────────
\begin{itemize}
	\item \textbf{Punkt 1:} Lorem ipsum dolor sit amet, consectetur adipiscing
		elit, sed do eiusmod tempor incididunt ut labore et dolore magna aliqua.
	\item \textbf{Punkt 2:} Ut enim ad minim veniam, quis nostrud exercitation
		ullamco laboris nisi ut aliquip ex ea commodo consequat.
	\item \textbf{Punkt 3:} Duis aute irure dolor in reprehenderit in voluptate
		velit esse cillum dolore eu fugat nulla pariatur.
\end{itemize}


\subsection{Problemstellung}
\label{sec:problemstellung}

Lorem ipsum dolor sit amet, consectetur adipiscing elit, sed do eiusmod tempor
incididunt ut labore et dolore magna aliqua. Ut enim ad minim veniam, quis nostrud
exercitation ullamco laboris nisi ut aliquip ex ea commodo consequat. Duis aute
irure dolor in reprehenderit in voluptate velit esse cillum dolore eu fugat nulla
pariatur.\autocite[Vgl.][S. 42]{muster:buch:2023}

% ── Beispiel: Aufzählung (nummeriert) ────────────────────────────────────────
\begin{enumerate}
	\item \textbf{Aspekt 1:} Lorem ipsum dolor sit amet, consectetur adipiscing
		elit, sed do eiusmod tempor incididunt ut labore et dolore magna aliqua.
	\item \textbf{Aspekt 2:} Ut enim ad minim veniam, quis nostrud exercitation
		ullamco laboris nisi ut aliquip ex ea commodo consequat.
	\item \textbf{Aspekt 3:} Duis aute irure dolor in reprehenderit in voluptate
		velit esse cillum dolore eu fugat nulla pariatur.
\end{enumerate}


\subsection{Zielsetzung und Forschungsfrage}
\label{sec:zielsetzung}

Lorem ipsum dolor sit amet, consectetur adipiscing elit, sed do eiusmod tempor
incididunt ut labore et dolore magna aliqua. Ut enim ad minim veniam, quis nostrud
exercitation ullamco laboris nisi ut aliquip ex ea commodo consequat.

% ── Beispiel: Hervorgehobener Textblock ───────────────────────────────────────
\begin{framed}
	\noindent\textbf{Forschungsfrage:} Lorem ipsum dolor sit amet, consectetur
	adipiscing elit, sed do eiusmod tempor incididunt ut labore et dolore magna
	aliqua?
\end{framed}

% ── Beispiel: Aufzählung (nummeriert) mit Verweis auf Unterkapitel ────────────
\begin{enumerate}
	\item \textbf{Teilziel 1:} Lorem ipsum dolor sit amet (s.\,Kapitel~\ref{sec:theorie}).
	\item \textbf{Teilziel 2:} Ut enim ad minim veniam (s.\,Kapitel~\ref{sec:methode}).
	\item \textbf{Teilziel 3:} Duis aute irure dolor (s.\,Kapitel~\ref{sec:hauptteil}).
	\item \textbf{Teilziel 4:} Excepteur sint occaecat (s.\,Kapitel~\ref{sec:fazit}).
\end{enumerate}


\subsection{Aufbau der Arbeit}
\label{sec:aufbau}

Kapitel~\ref{sec:theorie} erarbeitet die theoretischen Grundlagen. Darauf aufbauend
beschreibt Kapitel~\ref{sec:methode} das methodische Vorgehen. Kapitel~\ref{sec:hauptteil}
präsentiert die zentralen Ergebnisse. Abschließend fasst Kapitel~\ref{sec:fazit}
die Befunde zusammen und gibt einen Ausblick auf weiteren Forschungsbedarf.

\clearpage
